% ============================================================================
\section{From Language to Logic}
\label{sec:logic}
% Source: assignments/07_logic/index.html
% ============================================================================

Problems 1--5 are programming problems (formulas written with \texttt{logic.py}). We list the formulas in
mathematical notation so the report can be read on its own; the code and grader output are in the notebook. As the
handout notes, this assignment has no hidden tests: our formulas pass all 30 regular points and all 6 extra-credit
points.

% ----------------------------------------------------------------------------
\subsection{Propositional logic}
\codepart{7.1}
\begin{enumerate}[label=(1\alph*)]
  \item ``If it's summer and we're in California, then it doesn't rain.'' \quad
        $(\text{Summer}\wedge\text{California})\to\neg\text{Rain}$
  \item ``It's wet if and only if it is raining or the sprinklers are on.'' \quad
        $\text{Wet}\leftrightarrow(\text{Rain}\vee\text{Sprinklers})$
  \item ``Either it's day or night (but not both).'' \quad
        $(\text{Day}\vee\text{Night})\wedge\neg(\text{Day}\wedge\text{Night})$
\end{enumerate}

% ----------------------------------------------------------------------------
\subsection{First-order logic}
\codepart{7.2}
\begin{enumerate}[label=(2\alph*)]
  \item ``Every person has a parent.'' \quad
        $\forall x.\ \text{Person}(x)\to\exists y.\ \text{Parent}(x,y)$
  \item ``At least one person has no children.'' \quad
        $\exists x.\ \text{Person}(x)\wedge\neg\exists y.\ \text{Child}(x,y)$
  \item $\text{Father}$ in terms of $\text{Male}$ and $\text{Parent}$ ($\text{Father}(x,y)$: $x$ has father $y$). \quad
        $\forall x\,\forall y.\ \text{Father}(x,y)\leftrightarrow\bigl(\text{Male}(y)\wedge\text{Parent}(x,y)\bigr)$
  \item $\text{Granddaughter}$ in terms of $\text{Female}$ and $\text{Child}$. \\
        $\forall x\,\forall y.\ \text{Granddaughter}(x,y)\leftrightarrow
         \bigl(\text{Female}(y)\wedge\exists z.\ \text{Child}(x,z)\wedge\text{Child}(z,y)\bigr)$
\end{enumerate}
The definitions in 2c and 2d are biconditionals, because a definition must also rule out every other father or
granddaughter, not just license the intended ones.

% ----------------------------------------------------------------------------
\subsection{Liar puzzle}
\codepart{7.3a}
Suspects: John, Susan, Mark, Nicole. Mark: ``It wasn't me!'' John: ``It was Nicole!'' Nicole: ``It wasn't me, it
was Susan!'' Susan: ``Nicole's a liar.'' Exactly one person tells the truth, and exactly one person crashed the server.
With $T(x)$ for $\text{TellTruth}(x)$ and $S(x)$ for $\text{CrashedServer}(x)$, each statement is true exactly when its
speaker tells the truth:
\begin{align*}
&T(\text{mark})\leftrightarrow\neg S(\text{mark}) \qquad
 T(\text{john})\leftrightarrow S(\text{nicole}) \qquad
 T(\text{nicole})\leftrightarrow\bigl(\neg S(\text{nicole})\wedge S(\text{susan})\bigr)\\
&T(\text{susan})\leftrightarrow\neg T(\text{nicole}) \\
&\exists x.\ T(x)\wedge\forall y.\ \bigl(T(y)\to x=y\bigr) \qquad\qquad
 \exists x.\ S(x)\wedge\forall y.\ \bigl(S(y)\to x=y\bigr)
\end{align*}
Asking the knowledge base $S(\$x)$ (\texttt{grader.py 3a-run}) gives \textbf{Yes: mark}; No: john, nicole, susan.
\emph{Mark crashed the server}, and Susan is the only one telling the truth. Reasoning by hand: Susan and Nicole
cannot both lie (that would make Nicole truthful), so the single truth-teller is one of them. If it were Nicole,
then Susan crashed the server and Mark's ``it wasn't me'' would also be true, which gives two truth-tellers. So Susan
is truthful, Nicole and John lie (so neither Nicole nor Susan did it), and Mark lies, so it was Mark.

% ----------------------------------------------------------------------------
\subsection{Odd and even integers}
\codepart{7.4a}
With $\text{Succ}(x,y)$ for ``$x$'s successor is $y$'' and $\text{Larger}(x,y)$ for ``$x$ is larger than $y$'':
\begin{enumerate}[label=(\arabic*),start=0,nosep]
  \item $\forall x\,\exists y.\ \text{Succ}(x,y)\wedge\neg(x=y)\wedge\forall z.\ \bigl(\text{Succ}(x,z)\to y=z\bigr)$
  \item $\forall x.\ \bigl(\text{Even}(x)\vee\text{Odd}(x)\bigr)\wedge\neg\bigl(\text{Even}(x)\wedge\text{Odd}(x)\bigr)$
  \item $\forall x\,\forall y.\ \bigl(\text{Even}(x)\wedge\text{Succ}(x,y)\bigr)\to\text{Odd}(y)$
  \item $\forall x\,\forall y.\ \bigl(\text{Odd}(x)\wedge\text{Succ}(x,y)\bigr)\to\text{Even}(y)$
  \item $\forall x\,\forall y.\ \text{Succ}(x,y)\to\text{Larger}(y,x)$
  \item $\forall x\,\forall y\,\forall z.\ \bigl(\text{Larger}(x,y)\wedge\text{Larger}(y,z)\bigr)\to\text{Larger}(x,z)$
\end{enumerate}
Model checking (\texttt{grader.py 4a-run}) answers \textbf{Yes} to the query
$\forall x\,\exists y.\ \text{Even}(y)\wedge\text{Larger}(y,x)$: the theorem is entailed. Informally, if $x$ is odd
its successor is even and larger; if $x$ is even, the successor of its successor is even and, by transitivity,
larger than $x$.

% ----------------------------------------------------------------------------
\subsection{Semantic parsing (extra credit)}
\codepart{7.5}
The three grammar rules build the following logical forms (arguments are the matched words, predicates are
title-cased, so ``likes'' becomes $\text{Likes}$):
\begin{itemize}[nosep]
  \item \emph{every $N_1$ $V$ some $N_2$}: \quad $\forall x.\ N_1(x)\to\exists y.\ N_2(y)\wedge V(x,y)$
  \item \emph{there is some $N_1$ that every $N_2$ $V$}: \quad $\exists y.\ N_1(y)\wedge\forall x.\ N_2(x)\to V(x,y)$
  \item \emph{if a $N_1$ $V_1$ a $N_2$ then the former $V_2$ the latter}: \quad
        $\forall x\,\forall y.\ \bigl(N_1(x)\wedge N_2(y)\wedge V_1(x,y)\bigr)\to V_2(x,y)$
\end{itemize}
In the third rule the indefinite ``a'' inside a conditional is read universally (``any person who likes any
cat''), and ``the former''/``the latter'' are bound to the first and second variable. All 12 extra-credit tests
pass.

% ----------------------------------------------------------------------------
\subsection{Explainability of logic-based systems}

The four considerations: \emph{respect}; \emph{assessing the fairness of rules}; \emph{contesting and correcting
decisions}; \emph{ability to change user behaviour}. Context: the liar puzzle of Problem~3, where a conclusion could
have real consequences (e.g.\ someone being fired).

\problempart{6a}{What counts as an explanation}
\question{Besides the final answer, what does a general first-order logic system, like the one used in Problem 3,
return that could count as an explanation?}
\expect{1--2 sentences.}
\begin{solution}
It can return the \emph{proof}: the formulas it was given (the knowledge base) together with the chain of
inference steps (resolution steps or derivations) that lead from them to the conclusion. A model-checking system can
instead return the satisfying models, showing that in every world consistent with the facts Mark crashed the server
and in no such world anyone else did. Both make every rule used to reach the verdict inspectable.
\end{solution}

\problempart{6b}{Why it may not be adequate}
\question{Why might this explanation not be adequate? (Hint: model misspecification, or how the question is framed.)}
\expect{1--2 sentences.}
\begin{solution}
The proof is only as good as the premises: it shows that the conclusion follows \emph{from the facts we encoded},
not that those facts are true or complete. In the puzzle, the premise ``exactly one person tells the truth'' is an
unverified assumption, statements are forced into true/false although people can be mistaken rather than lying, and
the question presupposes that one of four named suspects is guilty. If any of this framing is wrong, the system
still produces a perfectly valid proof that blames the wrong person.
\end{solution}

\problempart{6c}{Which considerations are met}
\question{Are any of the four explainability considerations met by the explanations a logic-based system produces?}
\expect{2--4 sentences.}
\begin{solution}
Partly. \emph{Assessing the fairness of rules} and \emph{contesting and correcting decisions} are well served,
because every premise and inference step is explicit: an accused person can point to the specific premise they
dispute (``I was not lying, I was mistaken''), and removing or correcting it changes the conclusion transparently.
\emph{Respect} is only partly met: a formal proof is intelligible to experts but not automatically to the person
affected, so it has to be translated into plain language. \emph{Ability to change behaviour} is weakly served
for a one-off verdict like this, although for rule-based decisions in general (for example ``loans require
$X$'') explicit rules do tell people what they could change to obtain a different outcome.
\end{solution}

% ----------------------------------------------------------------------------
\subsection{Applications of soundness and completeness}

Sound: $\{f: KB\vdash f\}\subseteq\{f: KB\vDash f\}$. \quad Complete: $\{f: KB\vdash f\}\supseteq\{f: KB\vDash f\}$.

\problempart{7a}{Is an LLM sound or complete?}
\question{Treat whatever GPT-3.5 generates as what it ``derives''. Given that LLMs hallucinate, is GPT-3.5 sound,
complete, both or neither? Use examples and the definitions.}
\expect{2--3 sentences for each property (4--6 in total).}
\begin{solution}
\textbf{Not sound.} Soundness requires everything the system derives to be true (entailed by its knowledge).
Hallucinations are exactly derivations outside $\{f:KB\vDash f\}$, for example a confidently cited research paper that
does not exist or a wrong date for a historical event, so $\{f:KB\vdash f\}\not\subseteq\{f:KB\vDash f\}$.
\textbf{Not complete either.} Completeness requires that every true statement can be derived, but GPT-3.5 fails to
produce many entailed facts: it makes arithmetic errors on long multiplications, cannot know events after its
training cut-off, and can answer a question correctly in one phrasing and incorrectly in another. Hence
$\{f:KB\vDash f\}\not\subseteq\{f:KB\vdash f\}$, and GPT-3.5 is neither sound nor complete.
\end{solution}

\problempart{7b}{Safety-critical deployment}
\question{System A is sound but not complete, and System B is complete but not sound. Which would you deploy in a
safety-critical setting, where acting on a false conclusion has severe consequences, and why?}
\expect{2--3 sentences tied to the definitions.}
\begin{solution}
System A. Soundness guarantees $\{f:KB\vdash f\}\subseteq\{f:KB\vDash f\}$: everything it concludes is true, so
we never act on a false conclusion. Its incompleteness only means that it sometimes says ``I don't know'', which can be
handled safely (fall back to a conservative default or a human). System B may derive false statements, which is
exactly the failure mode the setting cannot tolerate.
\end{solution}

\problempart{7c}{With a separate verifier}
\question{With a separate verifier available (e.g.\ a human in the loop) and the goal of capturing as much true
information as possible, which system would you choose, and why?}
\expect{2--3 sentences tied to the definitions.}
\begin{solution}
System B. Completeness guarantees $\{f:KB\vDash f\}\subseteq\{f:KB\vdash f\}$: every true statement is among its
outputs, so nothing true is missed. Its unsoundness (some outputs are false) is acceptable because the verifier
filters the hypotheses, and the combination of B plus the verifier ends up both complete and, after filtering,
sound. System A would silently miss true facts that no verifier could recover.
\end{solution}
